\chapter{Conclusiones y Recomendaciones}

\section{Conclusiones}

El primer objetivo específico del proyecto planteó la integración de versiones adaptadas de los algoritmos de modulación y demodulación AM y ASK/OOK sobre una plataforma embebida de bajo costo. Este objetivo fue cumplido satisfactoriamente. Ambas técnicas fueron implementadas sobre el microcontrolador Raspberry Pi Pico utilizando MicroPython. Las pruebas realizadas confirmaron que ambas técnicas operan de forma continua y estable a lo largo de todo el rango disponible de sus parámetros de control, sin pérdida de sincronización ni degradación perceptible del flujo de procesamiento. El entregable asociado a este objetivo es el firmware funcional ejecutado sobre el RP2040, que integra generación, modulación y demodulación dentro de un ciclo continuo de procesamiento por bloques. Un aspecto central derivado de este proceso es que la adaptación al entorno embebido no consiste únicamente en trasladar formulaciones matemáticas a código ejecutable, sino en redefinir estructuras algorítmicas para conservar estabilidad operativa bajo restricciones reales de memoria, procesamiento y temporización.

Como segundo objetivo se contempló el desarrollo de un sistema de visualización y análisis que permitiera adquirir, reconstruir y representar continuamente las señales procesadas por el sistema embebido. Este objetivo fue alcanzado mediante el desarrollo de una aplicación de monitoreo basada en PyQt5, capaz de representar en tiempo real las tres etapas de la señal procesada. El entregable correspondiente es la aplicación de monitoreo funcional, que incorpora visualización gráfica continua, panel de parámetros operativos y registro de métricas de desempeño. La decisión central de este componente fue la estructuración de la transmisión mediante frames seriales sincronizados. Esta estrategia demostró ser determinante para preservar la alineación temporal entre señales durante la adquisición continua, eliminando la necesidad de mecanismos externos de reagrupamiento en la aplicación de monitoreo. Adicionalmente, el mecanismo de conmutación automática entre técnicas mediante metadata operó correctamente, demostrando así que la integración entre el firmware y la aplicación de monitoreo mantuvo una sincronización continua durante sesiones prolongadas de transmisión, confirmando la viabilidad de la arquitectura de comunicación propuesta.

Seguidamente, el tercer objetivo específico se orientó a la optimización del procesamiento y flujo de datos dentro del sistema embebido para garantizar un desempeño estable bajo condiciones de operación continua. Este objetivo fue cumplido mediante tres refinamientos concretos: la reserva estática de buffers desde el inicio de la ejecución, la externalización de la configuración operativa hacia un archivo (\texttt{config.json}) y la reorganización estructural del firmware bajo una interfaz genérica de módulos de procesamiento que desacopla la lógica interna de cada técnica del flujo global del sistema. La incorporación de estas estrategias permitió consolidar una arquitectura más estable, flexible y mantenible, reduciendo interrupciones durante la ejecución continua y facilitando la adaptación del sistema ante distintos escenarios de operación.

Por último, se planteó la evaluación de la calidad de los procesos de modulación y demodulación mediante métricas de precisión y estabilidad. Este objetivo fue alcanzado con la implementación de tres indicadores complementarios: la Correlación Cruzada Normalizada (NCC) para similitud estructural, el Error Cuadrático Medio Normalizado (NRMSE) para precisión numérica de reconstrucción y la desviación estándar del historial de NCC ($\sigma_{\text{NCC}}$) para estabilidad temporal del procesamiento continuo. El entregable correspondiente es el módulo de evaluación cuantitativa integrado al flujo de adquisición del monitor, con adaptación de la variante de NCC según la naturaleza analógica o digital de la técnica activa. La capacidad diferenciadora de los indicadores confirma que el sistema de evaluación implementado permite distinguir configuraciones con comportamiento visual similar y caracterizar objetivamente el efecto de variaciones de parámetros sobre el desempeño del sistema, cumpliendo adecuadamente su propósito dentro del alcance del proyecto.

En términos generales, el cumplimiento de estos objetivos demostró la viabilidad de implementar técnicas de modulación y demodulación sobre plataformas embebidas de bajo costo, integrando generación, procesamiento, transmisión y visualización de señales dentro de una solución funcional y estable. La solución responde directamente al problema planteado originalmente sobre disponer de una plataforma experimental funcional que permita estudiar y verificar el comportamiento de algoritmos de procesamiento digital de señales en condiciones reales de ejecución, sin depender de simuladores ni de equipamiento especializado. De esta manera, los objetivos específicos y entregables definidos fueron alcanzados satisfactoriamente, consolidando una base técnica orientada a la experimentación práctica y al análisis aplicado de técnicas de procesamiento digital de señales.

\subsection{Hallazgos técnicos relevantes}

El desarrollo del sistema evidenció que la ejecución interpretada de MicroPython sobre el RP2040 introduce restricciones operativas que van más allá de la capacidad de cómputo bruta del microcontrolador. La sobrecarga asociada a la interpretación dinámica de instrucciones, la manipulación continua de estructuras en memoria y las operaciones de empaquetado serial adquieren mayor relevancia conforme se incrementa la frecuencia de iteraciones o el tamaño del bloque procesado. Esta característica permitió evidenciar que la selección del entorno de ejecución influye directamente sobre la estabilidad y alcance funcional del sistema desarrollado, condicionando la forma en que deben estructurarse los ciclos de procesamiento dentro del entorno embebido.

Un hallazgo significativo desde la perspectiva del diseño es que la arquitectura modular desacoplada adoptada demostró reducir el ciclo de iteración durante el refinamiento del sistema. La separación entre lógica de procesamiento, configuración operativa y flujo de adquisición permitió modificar parámetros de operación sin intervenir constantemente sobre el firmware del microcontrolador, facilitando el proceso de experimentación y validación durante el desarrollo. Asimismo, esta organización contribuyó a mantener una estructura más flexible y mantenible, favoreciendo la integración de nuevas configuraciones sin comprometer el funcionamiento general del sistema.

La evaluación cuantitativa realizada permitió evidenciar que el comportamiento de las métricas implementadas depende directamente del enfoque de procesamiento por bloques utilizado en el sistema. Debido a que cada análisis se realiza sobre conjuntos discretos de muestras, variaciones de muy corta duración pueden no reflejarse completamente dentro de los indicadores calculados en cada iteración. No obstante, las métricas incorporadas demostraron ser suficientes para caracterizar de manera consistente la estabilidad y similitud de las señales reconstruidas durante las pruebas realizadas, permitiendo validar el comportamiento general del sistema dentro del alcance definido para el proyecto.

Finalmente, el proceso de desarrollo permitió comprobar que es posible integrar generación, procesamiento y transmisión de señales dentro de una misma plataforma embebida manteniendo un funcionamiento estable y continuo, evidenciando que la selección de estrategias adecuadas para el manejo y procesamiento de datos resulta fundamental para lograr que todos estos procesos operen de forma armoniosa dentro del mismo entorno, evitando sobrecarga computacional, pérdida de sincronización o interrupciones que comprometan la estabilidad general del sistema.

\section{Recomendaciones}

\subsection{Continuidad y expansión futura del proyecto}

La arquitectura modular presente en este proyecto constituye la base más directa para su expansión futura. Dado que cada técnica de modulación se encapsula como un módulo independiente que expone una interfaz genérica de procesamiento, la incorporación de nuevas técnicas no requiere modificar el flujo principal del sistema embebido ni la lógica de adquisición del monitor. Se recomienda extender el sistema hacia técnicas de modulación angular, particularmente FM y FSK, como siguiente paso natural de expansión. Estas técnicas introducen variaciones de frecuencia instantánea en lugar de amplitud, lo que exige redefinir la estrategia de generación de portadora dentro del entorno embebido y evaluar si el modelo de ejecución de MicroPython sobre el RP2040 puede sostener el cómputo de fase acumulada de forma continua sin comprometer la estabilidad temporal del sistema. Esta evaluación representaría en sí misma un aporte técnico relevante, dado que permitiría determinar los límites prácticos de la plataforma ante técnicas de mayor demanda computacional.

La evaluación de alternativas de implementación sobre la misma plataforma de hardware también representa una línea de continuidad importante para el proyecto. En particular, se recomienda analizar la migración del firmware hacia C/C++ sobre el RP2040 como alternativa al entorno interpretado de MicroPython. Esta transición es posible sin modificar la plataforma física y ayudaría a eliminar la sobrecarga asociada a la interpretación dinámica de instrucciones, ampliando la capacidad computacional disponible para incorporar técnicas de comunicación de mayor demanda. De igual manera, permitiría incorporar filtros de mayor orden para la etapa de demodulación, mejorar la selectividad frecuencial del procesamiento y reducir la latencia \textit{end-to-end} del sistema sin alterar la arquitectura general desarrollada. Alternativamente, se podría considerar migrar el sistema hacia plataformas con mayor capacidad de procesamiento, como la Raspberry Pi Pico 2 basada en el RP2350, entre otras, manteniendo el criterio de bajo costo que caracterizó la solución original.

La implementación desarrollada también deja abierta la posibilidad de profundizar en estrategias de paralelización utilizando la arquitectura dual-core del RP2040. La separación de responsabilidades entre ambos núcleos podría reducir la interferencia temporal entre tareas concurrentes y mejorar la uniformidad de la transmisión de frames. Este refinamiento no requeriría modificaciones sobre la aplicación de monitoreo ni sobre la arquitectura funcional general del sistema, por lo que representa una optimización de bajo costo de implementación con impacto potencial sobre la estabilidad temporal del procesamiento continuo.

\subsection{Aspectos no abordados y oportunidades de profundización}

El sistema de evaluación implementado caracteriza adecuadamente el desempeño bajo condiciones controladas de operación, pero no contempla la introducción deliberada de ruido, distorsión o pérdidas dentro del canal de procesamiento. Se recomienda incorporar como extensión futura un módulo de perturbación configurable que permita evaluar la robustez de las técnicas implementadas bajo condiciones adversas de señal. Esta capacidad permitiría analizar la degradación de los indicadores NCC y NRMSE en función del nivel de ruido introducido y establecer comparaciones más representativas entre AM y ASK/OOK bajo escenarios equivalentes de perturbación.

La evaluación realizada opera principalmente en el dominio del tiempo, por lo que existe una oportunidad clara de profundización mediante el análisis frecuencial de las señales procesadas. Se recomienda incorporar indicadores como coherencia espectral o relación señal-ruido espectral, complementando la información proporcionada actualmente por NCC y NRMSE. La integración de análisis espectral dentro del monitor, mediante transformada rápida de Fourier aplicada sobre los bloques recibidos, ampliaría considerablemente la capacidad diagnóstica del sistema sin requerir modificaciones sobre el firmware embebido.

En el caso específico de ASK/OOK, la variante digital de NCC implementada cuantifica coincidencia lógica entre estados binarios, pero no detecta errores de temporización sub-simbólicos ni desalineamientos parciales entre la secuencia original y la recuperada. Por esta razón, se recomienda estudiar la incorporación de una métrica de tasa de error de bit (BER) calculada sobre bloques suficientemente extensos para obtener representatividad estadística. Esta incorporación permitiría caracterizar el desempeño de ASK/OOK mediante indicadores más cercanos a los utilizados convencionalmente en la literatura de comunicaciones digitales y ampliaría el valor experimental del sistema.

Respecto a la aplicación de monitoreo, se recomienda incorporar un mecanismo de registro persistente capaz de almacenar el historial completo de métricas y parámetros operativos durante cada sesión de ejecución. Aunque la implementación actual calcula y presenta indicadores en tiempo real mediante una ventana deslizante, el sistema no conserva automáticamente los datos históricos una vez finalizada la ejecución. La posibilidad de exportar esta información en formatos estructurados facilitaría el análisis comparativo entre configuraciones y fortalecería la reproducibilidad experimental del sistema en contextos de laboratorio.

Por otro lado, se recomienda explorar la comunicación inalámbrica como alternativa a la transmisión serial por USB para futuras versiones del sistema. La incorporación de mecanismos inalámbricos de bajo costo permitiría eliminar la restricción de proximidad física entre el dispositivo embebido y la estación de monitoreo, ampliando las posibilidades de uso del sistema en contextos donde la movilidad sea relevante. Esta evolución implicaría adaptar la estructura de frames sincronizados, incorporando mecanismos de detección y corrección de errores adecuados a canales inalámbricos, donde la confiabilidad de transmisión no puede asumirse equivalente a la de una conexión USB directa.

Finalmente, el desarrollo de este proyecto permitió evidenciar la viabilidad de implementar técnicas de procesamiento digital de señales dentro de entornos computacionalmente limitados, considerando múltiples restricciones propias de este tipo de plataformas. La solución presentada consolidó una base funcional para el estudio, validación y experimentación de técnicas de modulación y demodulación, aportando además una arquitectura flexible con potencial de extensión hacia futuras aplicaciones y líneas de desarrollo dentro del área de sistemas embebidos.